%============================================================
% CH2 · PATTERN 4 — BACKTRACKING OVER PATHS: THEORY
%============================================================
\subsection{Pattern Theory}
\label{ssec:ch2-p4-theory}

\begin{patternbox}{Pattern 4 at a Glance}
Run DFS, but \textbf{remove the mark of a vertex when you leave it}. The marks now mean ``on the current path'' (the gray vertices), and the search walks through \textbf{every simple path} from the start --- one recursive call per path. Keep whatever the question needs (the path itself, its sum, its set of colours) \textbf{on the path}: add it when you enter a vertex, undo it when you leave. The cost is the number of simple paths, which is exponential: affordable for about $n \le 10$--$15$ vertices or $\le 25$ cells, or when pruning, a cap or a DAG keeps the count small.

\begin{center}
\begin{tikzpicture}[x=1cm, y=1cm]
    \begin{scope}
        \tikzset{g2-u1/.style=gvsrc, g2-u6/.style=gvdst}
        \pic {g2-G3};
        \node[font=\scriptsize\bfseries] at (2.15,-1.75) {$G_3$, start 1, target 6};
    \end{scope}
    \begin{scope}[xshift=6.3cm, yshift=1.3cm,
        tn/.style={rectangle, rounded corners=2pt, draw=darkgray!70, fill=white, font=\scriptsize\ttfamily,
                   inner sep=2pt, minimum height=4.2mm},
        tg/.style={tn, draw=accentorange, fill=answerfill, line width=0.9pt},
        tr/.style={tn, draw=successgreen, fill=basefill}]
        \node[tr] (n1) at (3.6,0) {1};
        \node[tn] (n12) at (1.5,-0.75) {1 2};
        \node[tn] (n14) at (5.7,-0.75) {1 4};
        \node[tn] (n123) at (0.2,-1.5) {1 2 3};
        \node[tn] (n125) at (2.4,-1.5) {1 2 5};
        \node[tn] (n145) at (5.7,-1.5) {1 4 5};
        \node[tn] (n1254) at (1.5,-2.25) {1 2 5 4};
        \node[tg] (n1256) at (3.3,-2.25) {1 2 5 6};
        \node[tn] (n1452) at (4.8,-2.25) {1 4 5 2};
        \node[tg] (n1456) at (6.6,-2.25) {1 4 5 6};
        \node[tn] (n14523) at (4.8,-3.0) {1 4 5 2 3};
        \foreach \a/\b in {n1/n12, n1/n14, n12/n123, n12/n125, n14/n145, n125/n1254, n125/n1256,
                           n145/n1452, n145/n1456, n1452/n14523}
            \draw[darkgray!80, -{Stealth[length=1.6mm]}] (\a) -- (\b);
        \node[note, text width=5.6cm, anchor=north] at (3.6,-3.35) {the search tree: one node per simple path from 1 (11 calls). Vertex 5 appears in two nodes --- marks are removed on exit};
    \end{scope}
\end{tikzpicture}
\end{center}
\begin{recurrencebox}
\centering\small
\texttt{go(u):}\ \ \textcolor{successgreen}{\textbf{enter}}: \texttt{onPath[u] = 1}, add $u$ to the state \quad record the path if it is an answer\\[2pt]
\texttt{for v in adj[u]: if !onPath[v] and promising(v): go(v)}\\[2pt]
\textcolor{accentorange}{\textbf{exit}}: \texttt{onPath[u] = 0}, undo exactly what \emph{enter} did\\[4pt]
time $=$ (number of search-tree nodes) $\times$ (degree) \qquad memory $=$ \bigO{n} (one path)
\end{recurrencebox}
\end{patternbox}

%------------------------------------------------------------
\subsubsection*{From the lessons to the pattern}

Tool~4 of Lesson~\ref{ssec:ch2-toolbox} made a single change to DFS: \texttt{onPath[u] = false} when $u$ is left. Plain DFS visits every \emph{vertex} once and is linear; backtracking visits every \emph{simple path} once and is exponential. The lesson's program printed the two paths $1, 2, 5, 6$ and $1, 4, 5, 6$ of $G_3$; the picture above shows everything it looked at on the way --- the \textbf{search tree}, whose nodes are the simple paths starting at 1 and whose edges are ``extend the path by one vertex''. The call stack at any moment is one root-to-node path of this tree, and the gray vertices are exactly the vertices of the current graph path.

Three things are new compared with the other patterns of this chapter:
\begin{itemize}[leftmargin=*]
    \item \textbf{counting the cost before coding} --- the search tree, not the graph, is what the program walks;
    \item \textbf{state on the path} --- a sum, a set of colours, a list of vertices --- maintained by ``add on enter, undo on exit'', the same discipline as the \texttt{onPath} marks themselves;
    \item \textbf{cutting the tree} --- pruning, a cap on the count, or memoising when two branches have the same future.
\end{itemize}

%------------------------------------------------------------
\subsubsection*{How many simple paths are there?}

\begin{center}
\small
\renewcommand{\arraystretch}{1.2}
\begin{tabular}{@{}lrrrrrr@{}}
\toprule
$k \times k$ grid graph, $k =$ & 1 & 2 & 3 & 4 & 5 & 6 \\
\midrule
paths corner $\to$ opposite corner & 1 & 2 & 12 & 184 & 8\,512 & 1\,262\,816 \\
all simple paths from a corner & 1 & 7 & 79 & 2\,111 & 153\,745 & 31\,811\,177 \\
\bottomrule
\end{tabular}

\smallskip
\begin{tabular}{@{}lrrrrrr@{}}
\toprule
complete graph $K_n$, $n =$ & 4 & 6 & 8 & 10 & 11 & 12 \\
\midrule
simple paths from one vertex & 16 & 326 & 13\,700 & 986\,410 & 9\,864\,101 & 108\,505\,112 \\
\bottomrule
\end{tabular}
\end{center}
Every number was counted by the backtracking program itself (the $6 \times 6$ row took 0.7\,s, $K_{12}$ took 2.3\,s). The complete graph is the worst case for $n$ vertices: a path from a fixed start is an ordered choice of distinct vertices, $\sum_{k} \frac{(n-1)!}{(n-1-k)!} \approx e\,(n-1)!$ paths. So:

\begin{diagrambox}
\begin{tikzpicture}[x=0.95cm, y=0.55cm]
    \draw[->, darkgray] (0,0) -- (0,9.2) node[above, font=\scriptsize]{$\log_{10}$(paths)};
    \foreach \y in {0,2,4,6,8} { \draw[darkgray] (-0.08,\y) -- (0.08,\y); \node[font=\tiny, anchor=east] at (-0.1,\y) {$10^{\y}$}; }
    \foreach \x/\v/\l in {1/1.08/3x3, 2/2.26/4x4, 3/3.93/5x5, 4/6.10/6x6}
        { \fill[primaryblue!60] (\x-0.3,0) rectangle (\x+0.3,\v); \node[font=\tiny, below] at (\x,0) {\l}; }
    \foreach \x/\v/\l in {6/1.20/$K_4$, 7/2.51/$K_6$, 8/4.14/$K_8$, 9/5.99/$K_{10}$, 10/8.04/$K_{12}$}
        { \fill[accentorange!70] (\x-0.3,0) rectangle (\x+0.3,\v); \node[font=\tiny, below] at (\x,0) {\l}; }
    \draw[warningred, dashed, thick] (0,8) -- (10.6,8) node[right, font=\scriptsize, align=left]{$\approx 10^8$ calls:\\2.3\,s here};
    \node[font=\scriptsize, text=primaryblue] at (2.5,7.0) {grid, corner to corner};
    \node[font=\scriptsize, text=accentorange] at (7.2,8.9) {$K_n$, from one vertex};
\end{tikzpicture}
\figcap{The two tables on a log scale. Each step of the grid size multiplies the count by 6 to 150; each extra vertex of $K_n$ multiplies it by $n$.}
\end{diagrambox}

%------------------------------------------------------------
\subsubsection*{When to recognise it (and when not)}

\textbf{Use backtracking when}
\begin{itemize}[leftmargin=*]
    \item the statement asks for \textbf{all} paths, or the best path under a rule that depends on the \emph{whole} path (``no vertex twice'', ``collect each cell's gold once'', ``a word along adjacent cells''), and
    \item the instance is \textbf{tiny}: $n \le 10$--$15$ vertices, at most $25$ cells of a grid, a word of $\le 15$ letters, a 7$\times$7 board with strong pruning; or
    \item the output itself has to be listed (then its size is the lower bound anyway: Problem~\ref{ssec:ch2-all-paths}); or
    \item the answer is \textbf{capped} (``print $\min(K, 10^6)$'') --- every call is one path, so the cap bounds the work (Exercise E4.4).
\end{itemize}
A special, linear case: on a \textbf{tree} there is exactly one path from the root to each vertex, so ``state along the path'' costs nothing extra --- the search tree \emph{is} the tree. The path-stack technique (template 3) is then a plain DFS.

\textbf{Do NOT use it when}
\begin{itemize}[leftmargin=*]
    \item only the \emph{number} of paths, or the best path, in a \textbf{DAG} is wanted: the future of a path depends only on its last vertex, so memoise --- DP on DAGs, \dpbook{12}{DP on DAGs \& Graphs};
    \item the shortest path is wanted: BFS (Chapter~3) or Dijkstra (Chapter~7);
    \item the question is ``visit every vertex'' (Hamiltonian, TSP) with $n \approx 15$--$20$: backtracking is $n!$, bitmask DP over (set, last vertex) is $2^n n^2$ --- \dpbook{13}{Bitmask DP} and Chapter~23;
    \item $n$ is $10^5$: no exponential search survives; look for a tree, a DAG or a DP.
\end{itemize}

%------------------------------------------------------------
\subsubsection*{The framework and its invariant}

\begin{center}
\small
\renewcommand{\arraystretch}{1.25}
\begin{tabularx}{\linewidth}{@{}>{\bfseries}l X@{}}
\toprule
Search tree & node $=$ simple path $s = x_0, x_1, \dots, x_k$; child $=$ the path extended by one neighbour of $x_k$ not on it \\
State & \texttt{onPath[]} plus whatever the answer needs: the vertex list, a running sum, colour counts, a bitmask \\
Enter $u$ & \texttt{onPath[u] = 1}; push $u$; add its contribution (sum $+{=}$ value, \texttt{cnt[colour]++}) \\
Exit $u$ & the exact inverse, in reverse order: \texttt{cnt[colour]--}, sum $-{=}$ value, pop, \texttt{onPath[u] = 0} \\
Answer & checked at every node (``best path anywhere'') or only at goal nodes (``paths to $t$'') \\
Cost & $\#\text{nodes} \times \deg$ time, \bigO{n} memory; the recursion depth is the path length $\le n$ \\
\bottomrule
\end{tabularx}
\end{center}

\begin{invariantbox}[Why backtracking sees every simple path exactly once]
\textbf{Invariant.} When \texttt{go(u)} is called, the call stack holds a simple path $x_0 = s, \dots, x_k = u$; \texttt{onPath} is true exactly for $x_0, \dots, x_k$, and every piece of path state equals its value for this path. When \texttt{go(u)} returns, \texttt{onPath} and the state are back to their values for $x_0, \dots, x_{k-1}$.

\emph{Proof} by induction on the call depth. Entering $u$ sets exactly the mark and the contributions of $u$; each recursive call returns with them unchanged (induction hypothesis); the exit undoes exactly the entry. So between two sibling calls the state is that of $x_0, \dots, x_k$.

\emph{Every path once.} A simple path $s = y_0, \dots, y_m$ is reached by the chain of calls $\texttt{go}(y_0) \to \dots \to \texttt{go}(y_m)$: each $y_{i+1}$ is a neighbour of $y_i$ that is not on the path $y_0..y_i$, so the loop of \texttt{go}$(y_i)$ calls it. It is reached only once, because the sequence of calls \emph{is} the path. Hence the search tree has one node per simple path from $s$. \hfill$\blacksquare$
\end{invariantbox}

%------------------------------------------------------------
\subsubsection*{State along the path: the path stack}

The marks are only the simplest piece of path state. Any quantity that is a ``sum'' of the vertices of the path can be kept the same way, \textbf{in \bigO{1} per step}, instead of being recomputed from the path at every node:

\begin{diagrambox}
\begin{tikzpicture}[x=0.95cm, y=0.95cm]
    \tikzset{gv/.append style={minimum size=6.8mm, font=\scriptsize\bfseries}}
    \node[gvseen] (r1) at (0,0) {1};
    \node[gvseen] (r4) at (1.6,-1) {4}; \node[gvdone] (r2) at (-1.6,-1) {2}; \node[gvdone] (r3) at (0,-1) {3};
    \node[gvdone] (r5) at (-2.3,-2) {5}; \node[gvdone] (r6) at (-0.9,-2) {6};
    \node[gvseen] (r7) at (1.6,-2) {7}; \node[gvcur] (r8) at (1.6,-3) {8};
    \foreach \a/\b in {1/2, 1/3, 2/5, 2/6} \draw[ge] (r\a) -- (r\b);
    \foreach \a/\b in {1/4, 4/7, 7/8} \draw[gtree] (r\a) -- (r\b);
    \foreach \v/\c in {1/a, 2/b, 3/a, 4/c, 5/b, 6/c, 7/a, 8/b}
        \node[font=\tiny\bfseries, text=softpurple, right=0pt of r\v] {\c};
    \node[note, text width=3.6cm] at (-1.6,-3.1) {black: entered and left again --- their colours were removed};
    % the stack
    \begin{scope}[xshift=3.7cm]
        \node[font=\scriptsize\bfseries] at (0.6,0.55) {call stack};
        \foreach \v/\c/\y in {1/a/0, 4/c/-0.75, 7/a/-1.5, 8/b/-2.25}
            \node[draw=memoorange, fill=hitfill, minimum width=1.3cm, minimum height=6.5mm, font=\scriptsize] at (0.6,\y) {\v\ (\c)};
        \draw[flow] (1.5,-2.25) -- (2.3,-2.25) node[right, font=\scriptsize, align=left]{enter 8: \texttt{cnt[b]++}};
        \draw[flow, accentorange] (2.3,-2.85) -- (1.5,-2.85) node[pos=0, right, font=\scriptsize, align=left]{exit 8: \texttt{cnt[b]--}};
    \end{scope}
    \begin{scope}[xshift=8.5cm]
        \node[font=\scriptsize\bfseries] at (0.8,0.55) {counts on the path 1, 4, 7, 8};
        \foreach \c/\k/\x in {a/2/0, b/1/0.8, c/1/1.6}{
            \node[cell] at (\x,-0.3) {\c};
            \node[curcell] at (\x,-1.0) {\k};
        }
        \node[font=\scriptsize, align=left, anchor=north west] at (-0.4,-1.6) {distinct colours $= 3$\\(the program prints 3 for 8)};
    \end{scope}
\end{tikzpicture}
\figcap{The tree $T$ of Lesson~\ref{ssec:ch2-toolbox} with colours a/b/c (template 3 uses $1,2,3$). While the DFS sits at 8, the stack is the root path $1, 4, 7, 8$ and \texttt{cnt[]} describes exactly that path.}
\end{diagrambox}

On a tree this is a \emph{linear} algorithm: every vertex is entered once and left once, so a question like ``for every vertex, does its colour already occur on the path from the root?'' (Exercise E4.5) or ``for every vertex, binary-search the prefix sums of its root path'' (Exercise E4.6) costs \VE{} plus the work per step. In a general graph the same discipline runs inside the exponential search; the state then has to be undone just as carefully as the marks.

%------------------------------------------------------------
\subsubsection*{Cutting the search tree}

Pruning means: at a node, prove that no extension can matter, and return. It is only allowed when the proof is airtight --- a wrong prune loses answers silently.

\begin{diagrambox}
\begin{tikzpicture}[x=0.62cm, y=0.62cm]
    \tikzset{c/.style={gcell, minimum size=6mm, font=\scriptsize}, v/.style={c, fill=hitfill}, cu/.style={c, fill=currentfill, draw=accentorange, line width=0.9pt},
             f/.style={c, fill=basefill}}
    % (a) bound
    \begin{scope}
        \node[font=\scriptsize\bfseries, align=center] at (1.5,1.05) {(a) bound};
        \foreach \x/\y/\s/\t in {0/0/v/2, 1/0/v/3, 2/0/cu/1, 3/0/c/1, 0/-1/c/1, 1/-1/c/2, 2/-1/c/1, 3/-1/c/1}
            \node[\s] at (\x,\y) {\t};
        \node[font=\scriptsize, align=center, text width=3.4cm] at (1.5,-2.5) {sum 6, unvisited 6;\\best so far 12 $\ge 6 + 6$:\\cut};
    \end{scope}
    % (b) split
    \begin{scope}[xshift=4.9cm]
        \node[font=\scriptsize\bfseries] at (1.5,1.05) {(b) the free cells split};
        \foreach \x in {0,1,2,3} \node[v] at (\x,0) {};
        \node[v] at (3,-1) {}; \node[cu] at (3,-2) {};
        \node[f] at (2,-2) {L}; \node[f] at (4,-2) {R};
        \node[gwall, minimum size=6mm] at (3,-3) {};
        \foreach \x/\y in {0/-1, 1/-1, 2/-1, 4/-1, 4/0, 0/-2, 1/-2} \node[c] at (\x,\y) {};
        \draw[flow, accentorange] (3,-1.35) -- (3,-1.65);
        \node[font=\scriptsize, align=center, text width=3.6cm] at (2,-4.75) {moving down, the wall is ahead, L and R are both free: the path can reach only one side};
    \end{scope}
    % (c) parity
    \begin{scope}[xshift=10.2cm]
        \node[font=\scriptsize\bfseries] at (1.5,1.05) {(c) parity};
        \foreach \x in {0,...,3} \foreach \y in {0,-1,-2} {
            \pgfmathtruncatemacro{\p}{mod(\x-\y+10,2)}
            \ifnum\p=0 \node[c, fill=primaryblue!12] at (\x,\y) {}; \else \node[c] at (\x,\y) {}; \fi }
        \node[cu] at (0,0) {s}; \node[c, fill=answerfill, draw=accentorange, line width=0.9pt] at (3,-2) {t};
        \node[font=\scriptsize, align=center, text width=3.6cm] at (1.5,-3.75) {every step changes the colour: a walk $s \to t$ has odd length here, so ``exactly 6 steps'' is impossible};
    \end{scope}
\end{tikzpicture}
\figcap{Three proofs that cut a branch. (a) An optimistic \textbf{bound} --- current value plus everything still available --- cannot beat the best: template 2. (b) A \textbf{dead-end test} --- some free cell can no longer be reached by the rest of the path: Exercise E4.7. (c) A \textbf{parity} argument on the bipartite grid.}
\end{diagrambox}

\begin{itemize}[leftmargin=*]
    \item \textbf{Bounds.} ``value so far $+$ optimistic value of the rest $\le$ best'' $\Rightarrow$ return. The optimistic rest can be crude (the sum of all unvisited values) or sharper (only the values still \emph{reachable} from the current vertex --- Problem~\ref{ssec:ch2-max-gold} measures the difference). Once the best is equal to the global maximum possible, stop the whole search.
    \item \textbf{Dead ends.} The target is no longer reachable from $u$ without crossing the path; or a cell that must still be visited can no longer be reached; or a cell has only one free neighbour left and therefore must be entered \emph{now}. On a DAG, ``can $v$ reach $t$ at all?'' is one reverse search (Pattern~5) and removes whole subgraphs before the search starts.
    \item \textbf{Parity and counting.} In a grid every step flips the colour of a chessboard; a path of exactly $L$ steps from $s$ to $t$ needs $L \equiv \text{dist}(s,t) \pmod 2$. Similar counting arguments: ``the remaining letters of the word occur often enough on the board'' (Exercise E4.2).
    \item \textbf{Order.} Try the most promising neighbour first (biggest value, rarest letter), so that a good ``best'' is known early and the bound bites sooner.
    \item \textbf{Caps.} If the answer is $\min(K, \text{cap})$, stop when the counter reaches the cap. Since every call is a new path, the number of calls is at most the cap.
\end{itemize}

%------------------------------------------------------------
\subsubsection*{When the future does not depend on the path: memoise}

Backtracking pays for the \emph{past}: two different paths that end in the same vertex are two different nodes of the search tree, because their sets of used vertices differ. In a \textbf{DAG} the past does not matter --- a path can never come back to a vertex it has used, so the possible continuations of a path depend only on its last vertex. The search tree then collapses into the DAG itself:

\begin{diagrambox}
\begin{tikzpicture}[x=1.05cm, y=0.9cm]
    \begin{scope}
        \node[gvsrc] (a0) at (0,0) {0}; \node[gv] (a1) at (1.4,0.8) {1}; \node[gv] (a2) at (1.4,-0.8) {2};
        \node[gvcur] (a3) at (2.8,0) {3}; \node[gv] (a4) at (4.2,0.8) {4}; \node[gvdst] (a5) at (4.2,-0.8) {5};
        \foreach \a/\b in {0/1, 0/2, 1/3, 2/3, 3/4, 3/5, 4/5} \draw[gde] (a\a) -- (a\b);
        \node[font=\scriptsize\bfseries] at (2.1,-1.6) {a DAG: 4 paths $0 \leadsto 5$};
    \end{scope}
    \begin{scope}[xshift=6.0cm, yshift=0.9cm, tn/.style={rectangle, rounded corners=2pt, draw=darkgray!70, fill=white,
                  font=\tiny\ttfamily, inner sep=1.5pt}, tz/.style={tn, fill=currentfill, draw=accentorange}]
        \node[tn] (s0) at (1.6,0) {0};
        \node[tn] (s1) at (0.7,-0.6) {01}; \node[tn] (s2) at (2.5,-0.6) {02};
        \node[tz] (s13) at (0.7,-1.2) {013}; \node[tz] (s23) at (2.5,-1.2) {023};
        \node[tn] (s134) at (0.3,-1.8) {0134}; \node[tn] (s135) at (1.1,-1.8) {0135};
        \node[tn] (s234) at (2.1,-1.8) {0234}; \node[tn] (s235) at (2.9,-1.8) {0235};
        \node[tn] (s1345) at (0.3,-2.4) {01345}; \node[tn] (s2345) at (2.1,-2.4) {02345};
        \foreach \a/\b in {s0/s1, s0/s2, s1/s13, s2/s23, s13/s134, s13/s135, s23/s234, s23/s235, s134/s1345, s234/s2345}
            \draw[darkgray!70] (\a) -- (\b);
        \draw[decorate, decoration={brace, amplitude=4pt, mirror}, accentorange] (-0.2,-2.7) -- (3.4,-2.7)
            node[midway, below=4pt, font=\scriptsize, text=accentorange, align=center]{the two subtrees under 3 are identical};
    \end{scope}
    \begin{scope}[xshift=11.0cm]
        \node[font=\scriptsize\bfseries] at (1.1,1.35) {memo: paths from $v$ to 5};
        \foreach \v/\c/\x/\y in {0/4/0/0, 1/2/0.9/0.6, 2/2/0.9/-0.6, 3/2/1.8/0, 4/1/2.7/0.6, 5/1/2.7/-0.6}{
            \node[gv, minimum size=6mm, font=\scriptsize\bfseries] (m\v) at (\x,\y) {\v};
            \node[gdist, above=0pt of m\v] {\c};
        }
        \foreach \a/\b in {0/1, 0/2, 1/3, 2/3, 3/4, 3/5, 4/5} \draw[gde] (m\a) -- (m\b);
        \node[font=\scriptsize, align=center] at (1.35,-1.55) {$P(v) = \sum_{v \to w} P(w)$, $P(5) = 1$};
    \end{scope}
\end{tikzpicture}
\figcap{Left: a DAG. Middle: its search tree from 0 --- vertex 3 is reached by two paths and everything below it is computed twice. Right: memoising on the last vertex merges the copies; the count is a DP over the DAG in \VE.}
\end{diagrambox}

\begin{itemize}[leftmargin=*]
    \item \textbf{Count or optimise on a DAG:} memoise on the last vertex --- \VE{}, even when the number of paths is astronomically large (\dpbook{12}{DP on DAGs \& Graphs}). Listing the paths still costs their total length; Problem~\ref{ssec:ch2-all-paths} compares both.
    \item \textbf{General graph, ``which vertices are used'' matters:} the honest state is (set of used vertices, last vertex) --- $2^n \cdot n$ states. For $n \le 20$ this bitmask DP beats $n!$ backtracking by far (\dpbook{13}{Bitmask DP}); for a grid with 25 gold cells it does not help, because almost every (set, last) pair is reached by only one path (Problem~\ref{ssec:ch2-max-gold}, Approach~3).
    \item \textbf{Trees:} one path per vertex --- nothing to memoise, the path stack is already linear.
\end{itemize}

%------------------------------------------------------------
\subsubsection*{Variations map}

\begin{diagrambox}
\begin{tikzpicture}[
    var/.style={rectangle, rounded corners=2pt, draw=patternborder, fill=patternbg,
                font=\scriptsize, align=center, text width=3.25cm, minimum height=1.3cm, inner sep=3pt}]
    \node[var, fill=answerfill, draw=accentorange, text width=5.6cm, font=\scriptsize\bfseries] (root) at (0,0)
        {un-mark on exit: every simple path; state added on enter, undone on exit};
    \node[var] (v1) at (-5.4,-2) {\textbf{list all paths $s \to t$}\\in a DAG, no marks needed\\(All Paths, 2.4.1)};
    \node[var] (v2) at (-1.8,-2) {\textbf{best path, any start}\\grid, temporary zeroing\\(Max Gold, 2.4.2)};
    \node[var] (v3) at (1.8,-2)  {\textbf{count paths through all}\\Hamiltonian, $n \le 8$\\(E4.1)};
    \node[var] (v4) at (5.4,-2)  {\textbf{a word along a path}\\letter-by-letter pruning\\(E4.2)};
    \node[var] (v5) at (-5.4,-3.6) {\textbf{max-weight simple path}\\$n \le 10$, every start\\(E4.3)};
    \node[var] (v6) at (-1.8,-3.6) {\textbf{count with a cap}\\$\min(K, 10^6)$, $n = 2\cdot10^5$\\(E4.4)};
    \node[var] (v7) at (1.8,-3.6)  {\textbf{path stack on a tree}\\colour counts, prefix sums\\(E4.5, E4.6)};
    \node[var] (v8) at (5.4,-3.6)  {\textbf{heavy pruning}\\$7 \times 7$ grid, 48 steps\\(E4.7)};
    \node[draw=patternborder, dashed, rounded corners=4pt, fit=(v1)(v4)(v5)(v8), inner sep=5pt] (grp) {};
    \draw[flow] (root.south) -- (grp.north);
    \node[note, below=2mm of grp] {count or optimise on a DAG: DP (DP book Ch.~12) $\cdot$ ``visit all'' with $n \approx 20$: bitmask DP (DP book Ch.~13, Chapter~23)};
\end{tikzpicture}
\end{diagrambox}

\textbf{Twists and how to handle them}
\begin{itemize}[leftmargin=*]
    \item \textbf{Any start, any end:} loop over every start vertex; record the answer at \emph{every} node, not only at a target.
    \item \textbf{Grid instead of graph:} mark a used cell by overwriting it (\texttt{0}, \texttt{'\#'}) and restore the old value on exit --- no separate \texttt{onPath} array (Problem~\ref{ssec:ch2-max-gold}).
    \item \textbf{Stop at the target:} ``paths from $s$ to $t$'' must not continue \emph{through} $t$; a path that reaches $t$ early may also be a failure (``exactly 48 steps'').
    \item \textbf{Undirected graph:} a path and its reverse are the same path; if they must be counted once, count ordered paths and divide by 2, or fix an orientation (e.g.\ start $<$ end).
    \item \textbf{Deep recursion:} the depth is only the path length, $\le n$ --- harmless for $n \le 25$, but a capped count on $2\cdot10^5$ vertices can build a path of $2\cdot10^5$ vertices: write it iteratively and reset the iterator of a vertex when it is left (Exercise E4.4).
\end{itemize}

%------------------------------------------------------------
\subsubsection*{Code templates}

\begin{lstlisting}[style=cpp]
#include <bits/stdc++.h>
using namespace std;
typedef long long ll;

// (1) ALL SIMPLE PATHS s -> t (any graph). onPath[] = the gray vertices = the current path.
//     limit: stop after that many paths (the answer may be exponential).
struct AllPaths {
    const vector<vector<int>>& adj;
    vector<char> onPath;
    vector<int> path;
    vector<vector<int>> out;
    int t;
    size_t limit;
    AllPaths(const vector<vector<int>>& g) : adj(g), onPath(g.size(), 0) {}
    void go(int u) {
        if (out.size() >= limit) return;              // cap reached: stop the whole search
        onPath[u] = 1; path.push_back(u);             // enter: extend the path
        if (u == t) out.push_back(path);              // a simple path ends here (do not continue past t)
        else
            for (int v : adj[u])
                if (!onPath[v]) go(v);
        onPath[u] = 0; path.pop_back();               // exit: undo exactly what enter did
    }
    vector<vector<int>> run(int s, int target, size_t lim) {
        t = target; limit = lim; out.clear();
        go(s);
        return out;
    }
};

// (2) BEST SIMPLE PATH FROM ANY START, vertex values val[v] >= 0, with an optimistic bound:
//     if (sum so far) + (sum of all values not on the path) <= best, no extension can win.
struct BestPath {
    const vector<vector<int>>& adj;
    const vector<int>& val;
    vector<char> onPath;
    ll best = 0, rest = 0, calls = 0;                 // rest = total value of vertices NOT on the path
    BestPath(const vector<vector<int>>& g, const vector<int>& w) : adj(g), val(w), onPath(g.size(), 0) {}
    void go(int u, ll sum) {
        calls++;
        onPath[u] = 1; sum += val[u]; rest -= val[u]; // enter
        best = max(best, sum);
        if (sum + rest > best)                        // prune: even taking everything left cannot beat best
            for (int v : adj[u])
                if (!onPath[v]) go(v, sum);
        onPath[u] = 0; rest += val[u];                // exit
    }
    ll run() {
        for (int x : val) rest += x;
        for (int s = 0; s < (int)adj.size(); s++) go(s, 0);
        return best;
    }
};

// (3) PATH-STACK DFS ON A TREE: the root path is the state. cnt[c] = how often colour c is on the
//     path root..u; add at enter, remove at exit. Iterative (trees can be 1e5+ deep).
//     Output: distinct[u] = number of different colours on the path root..u.
vector<int> distinctOnRootPath(const vector<vector<int>>& adj, const vector<int>& col, int root) {
    int n = adj.size(), cur = 0;
    vector<int> cnt(*max_element(col.begin(), col.end()) + 1, 0), distinct(n, 0), par(n, -1);
    vector<size_t> it(n, 0);
    vector<int> st;
    st.push_back(root); par[root] = root;
    if (cnt[col[root]]++ == 0) cur++;                 // enter root
    distinct[root] = cur;
    while (!st.empty()) {
        int u = st.back();
        if (it[u] < adj[u].size()) {
            int v = adj[u][it[u]++];
            if (v == par[u]) continue;
            par[v] = u;
            if (cnt[col[v]]++ == 0) cur++;            // enter v: its colour joins the path
            distinct[v] = cur;
            st.push_back(v);
        } else {
            if (--cnt[col[u]] == 0) cur--;            // exit u: its colour leaves the path
            st.pop_back();
        }
    }
    return distinct;
}

int main() {
    // G3 of the lessons, 0-indexed: edges 1-2 1-4 2-3 2-5 4-5 5-6 (vertex 7 isolated)
    vector<vector<int>> G3(7);
    int E[6][2] = {{1, 2}, {1, 4}, {2, 3}, {2, 5}, {4, 5}, {5, 6}};
    for (auto& e : E) { G3[e[0] - 1].push_back(e[1] - 1); G3[e[1] - 1].push_back(e[0] - 1); }
    AllPaths ap(G3);
    for (auto& p : ap.run(0, 5, 1000)) {               // 1 2 5 6 / 1 4 5 6
        for (int x : p) cout << x + 1 << " ";
        cout << "\n";
    }
    cout << ap.run(0, 5, 1).size() << "\n";            // 1 (capped)

    vector<int> val = {3, 1, 4, 1, 5, 9, 0};           // G3 with vertex values (7 is isolated)
    BestPath bp(G3, val);
    ll best = bp.run();
    cout << best << " " << bp.calls << "\n";           // 23 32   (63 calls without the bound)

    // the tree T of Lesson 5, 0-indexed, colours of 1..8
    vector<vector<int>> T(8);
    int F[7][2] = {{1, 2}, {1, 3}, {1, 4}, {2, 5}, {2, 6}, {4, 7}, {7, 8}};
    for (auto& e : F) { T[e[0] - 1].push_back(e[1] - 1); T[e[1] - 1].push_back(e[0] - 1); }
    vector<int> col = {1, 2, 1, 3, 2, 3, 1, 2};
    for (int d : distinctOnRootPath(T, col, 0)) cout << d << " ";   // 1 2 1 2 2 3 2 3
    cout << "\n";
}
\end{lstlisting}

On the values $3, 1, 4, 1, 5, 9$ of $G_3$ the best path is $3, 2, 1, 4, 5, 6$ with value 23, which uses the whole component; once it is found, $\text{sum} + \text{rest} = 23$ never beats it again, and the bound halves the search from 63 to 32 calls.

\begin{insightbox}[Which template when?]
\renewcommand{\arraystretch}{1.25}
\begin{tabularx}{\linewidth}{@{}>{\bfseries}l X@{}}
(1) \texttt{AllPaths} & list or count the simple paths between two vertices; with a \texttt{limit} when the count is capped; in a DAG the \texttt{onPath} test can be dropped (Problem~\ref{ssec:ch2-all-paths}) \\
(2) \texttt{BestPath} & the best simple path anywhere (max value, max length); every node is a candidate; add the strongest bound you can prove \\
(3) \texttt{distinctOnRootPath} & any ``for every vertex of a tree, something about its root path'' --- linear, iterative; replace the colour counts by a sum, a stack of prefix sums, a set \\
\end{tabularx}
\end{insightbox}

%------------------------------------------------------------
\subsubsection*{Common mistakes}

\begin{warningbox}[Mistake 1 --- forgetting to un-mark]
\begin{lstlisting}[style=snippet]
void go(int u) { onPath[u] = true; for (int v : adj[u]) if (!onPath[v]) go(v); }        // plain DFS
void go(int u) { onPath[u] = true; for (int v : adj[u]) if (!onPath[v]) go(v); onPath[u] = false; }
\end{lstlisting}
The first version is a correct DFS and a wrong backtracking: on $G_3$ it prints only $1, 2, 5, 6$ --- vertex 5 stays marked and the path $1, 4, 5, 6$ is never tried.
\end{warningbox}

\begin{warningbox}[Mistake 2 --- undoing on only one exit]
An early \texttt{return} (``target reached'', ``pruned'') that skips the un-mark leaves a vertex gray forever. Put the undo at the single end of the function, or undo before every \texttt{return}. In a grid that is marked by overwriting, restoring the wrong value (\texttt{1} instead of the saved gold) is the same bug.
\end{warningbox}

\begin{warningbox}[Mistake 3 --- continuing through the target]
For ``paths from $s$ to $t$'' a path that reaches $t$ is finished. If the search goes on, it also counts walks like $s \dots t \dots t'$ that pass through $t$ --- in a DAG this does no harm only if $t$ is a sink.
\end{warningbox}

\begin{warningbox}[Mistake 4 --- backtracking where a DP is needed]
``Count the paths from 1 to $n$ in a DAG with $n = 10^5$'' has up to $2^{n-2}$ paths --- no search can list them. The future depends only on the last vertex: memoise (and take the count modulo the given prime).
\end{warningbox}

\begin{warningbox}[Mistake 5 --- a prune that is only ``usually'' right]
``Start only at cells with few neighbours'', ``never turn twice in a row'' --- such rules can silently drop the optimum. Prove them, or test them against the unpruned search on thousands of random inputs before trusting them; Problem~\ref{ssec:ch2-max-gold} shows a rule that passed every test we could think of but has no proof, and is therefore kept optional.
\end{warningbox}

\begin{checklistbox}[``I'm stuck'' checklist for Pattern 4]
\begin{itemize}[leftmargin=*]
    \item[$\square$] How many simple paths can there be in the worst case? Count them (by the formula, or by running the search on the worst input) before coding.
    \item[$\square$] Is the graph a DAG or a tree? Then the past does not matter: no marks needed (DAG) or one path per vertex (tree) --- maybe memoise instead.
    \item[$\square$] What does the answer depend on --- only the last vertex, or the set of used vertices? That decides between DP on the last vertex, bitmask DP and backtracking.
    \item[$\square$] Is every piece of state undone on every way out of the function?
    \item[$\square$] Which bound, dead-end test or parity argument can cut a branch, and can I prove it?
    \item[$\square$] Is the answer capped, or does the search stop at the first success? Then stop the whole recursion, not just the current call.
    \item[$\square$] Can the path be $10^5$ long? Then write the search iteratively.
\end{itemize}
\end{checklistbox}
